\section{The 481-phase application and a fixed error budget}\label{sec:numerics}

\subsection{One concrete finite input}
We use the field and reverse ranks defined in Section~\ref{sec:certificates},
with spatial grid \(G=256\), time scale \(249/250\), and 481 nonempty phases.
There are 242,176 positive type/cells; each is used exactly once. By
Lemma~\ref{lem:reverse}, their reverse order defines masks in \([0,1]\),
and the checked background field is the cumulative field at the end of
the corresponding forward phase. The four types use a common endpoint.
This construction replaces the earlier candidate input; its mass is a
direct sum over disjoint type/cells, with no unverified union-area correction.

The exact mass and duration are
\begin{equation}\label{eq:input481numbers}
 A=\frac{3842853135627}{2560000000000}=1.501114506104296875,
 \qquad T=\frac{30088384743729}{25000000000}.
\end{equation}
The field is already time-scaled in these formulas. No further factor
\((249/250)^2\) is applied to its squared-chord risk. The coefficient
of the twelve integrals is \(D_*/24=19/750\).

\begin{center}
\begin{tabular}{@{}ll@{}}
\toprule
Quantity & Exact value or choice\\
\midrule
Target \(c\) & \(1501/1000\)\\
Designated clock cap \(K\) & \(255/128\)\\
Ordinary capacity cap \(\lambda\) & \(4/5\)\\
Ordinary credit floor \(\rho\) & \(1/8\)\\
Continuum risk bound \(r_*\) & \(9999/10000\)\\
Weight exponent \(d\) & \(19999/20000\)\\
Virtual exponent \(\gamma\) & \(39999/40000\)\\
Geometric and correction exponents \(\xi,\zeta\) & \(10^{-6},10^{-6}\)\\
Finite-sieve coefficient \(D_*\) & \(76/125\)\\
Spatial grid; designated direction count & \(256;\ 16\)\\
Total early-retirement probability budget & \(10^{-5}\)\\
\bottomrule
\end{tabular}
\end{center}

The capacity bounds for this same field are
\begin{align*}
 U_H&=\frac{159305594247}{80000000000},
 &K-U_H&=\frac{69405753}{80000000000},\\
 U_L&=\frac{2535641192953}{3200000000000},
 &\lambda-U_L&=\frac{24358807047}{3200000000000},\\
 U_{\rm tail}&=\frac{294190968417}{400000000000},
 &1-\rho-\lambda&=\frac3{40}.
\end{align*}
All margins are strict. The finite calculation covers 144 direction/type-pair
rows, including both axes; six tail bounds cover all heights at least eight.

The maximum in the freshly recomputed complete angular-replay records is
\begin{equation}\label{eq:recordedrisk481}
 r_{\rm rec}=\frac{93740574637267}{93750000000000},\qquad
 r_*-r_{\rm rec}=\frac{50362733}{93750000000000}>0.
\end{equation}
It occurs at reverse rank 205, type 2, microcell \((5,238)\).
Every positive type/cell is covered by exactly one phase record.
The checking algorithms and their soundness are described in
Section~\ref{sec:certificates}; the execution status for this edition
is stated separately in Section~\ref{sec:status}.

The finite sieve consists of the 175 odd primes from 3 to 1049, with
\[
 D_P=\frac34\prod_p(1-p^{-2})<\frac{76}{125}.
\]
For the discrete pair sum set
\begin{equation}\label{eq:adiscrete}
 a=r_*+\frac1{320000}=\frac{319969}{320000}<d.
\end{equation}
Equation~\eqref{eq:discrete481} gives
\(\mathcal B_n^{\Gamma^c}\le a\log n\) once
\(\log n\ge3\cdot10^{18}\). This explicitly pays the additive
lattice error; its coefficient is not silently left equal to \(r_*\).
The exponent conditions hold because
\[
 1-d-2\xi-\zeta=\frac{47}{10^6}>0,\qquad
 \gamma-d=\frac1{40000}>0.
\]

\subsection{Three global causes with fixed orders}
All moment orders are chosen before taking \(n\) to infinity.
For the pair-correction and virtual-deletion guards use
\(m_{\rm guard}=2^{18}=262144\). Including the factor \(B=n^d\)
for transport to the root point law, the powers in
\eqref{eq:etatail} and \eqref{eq:virtualtail} are respectively
\begin{align*}
 4-m+2m\xi+d(m-1)+m\zeta+d&=-8.320768,\\
 3-\gamma m+d(m-1)+\xi+d&=-3.553599.
\end{align*}
Each of these two probabilities is therefore less than \(2\cdot10^{-6}\),
uniformly in the root, for sufficiently large \(n\). Their constants
are finite for this fixed order; no unspecified constant is assigned
an assumed numerical value.

For designated-line clocks, \(C_{\rm line}=4G=1024\) and
\[
 \varepsilon_H=\frac{K-U_H}{8C_{\rm line}T}
 =\frac{115676255}{1314581588376682496}.
\]
Equations~\eqref{eq:clockexplicit}--\eqref{eq:clockthreshold} give
the same budget \(2\cdot10^{-6}\). In particular,
\(\log n\ge10^7\) suffices for this clock estimate. This is a
threshold for one cause, not a value of the final \(n_0\).

\subsection{The weight cap with a fixed moment order}
By \eqref{eq:hazarddecomp}, for every fixed integer \(m\), uniformly
in the root,
\[
 I_p\le Y_n+R_n,\qquad
 \|Y_n\|_{L^m(\Prob^{(p)})}\le(r_*+o_m(1))\log n,
 \qquad \E^{(p)}R_n=o(\log n),
\]
where \(Y_n,R_n\ge0\). Put
\[
 h=\frac{d-r_*}{4}=\frac1{80000},\qquad
 b=\frac{d+r_*}{2}=\frac{39997}{40000},\qquad
 m_{\rm cap}=2^{22}=4194304.
\]
After increasing \(n\) by a root-independent threshold, we have
\begin{equation}\label{eq:fixedriskmoment}
 \|Y_n\|_{m_{\rm cap}}\le(r_*+h/2)\log n,\qquad
 \E^{(p)}R_n\le10^{-6}h\log n.
\end{equation}
For \(q=(r_*+h/2)/(b-h)=159985/159986\), Bernoulli's inequality
gives \(q^{-159985}\ge2\). Since \(26\cdot159985\le m_{\rm cap}\),
Markov's inequality yields
\begin{equation}\label{eq:fixedcapbudget}
 \Prob^{(p)}(I_p>b\log n)
 \le q^{m_{\rm cap}}+10^{-6}
 \le2^{-26}+10^{-6}<2\cdot10^{-6}.
\end{equation}
The two terms pay for \(Y_n>(b-h)\log n\) and \(R_n>h\log n\)
separately.

Both axes are designated directions, so an active root has
\(\sum_{\ell\ni p,\,\ell\in\Gamma} b_{N_\ell}\ge1\) in
Lemma~\ref{lem:product}. Its self-correction
contribution is nonpositive. Each designated factor is at most
\(2/(2-K)=256\), and there are 16 designated directions. Thus
\[
 \log w_p\le I_p+16\log256<I_p+128.
\]
Reaching \(w_p=n^d\) requires \(I_p>d\log n-128\).
If \(\log n>5120000\), this implies the event in
\eqref{eq:fixedcapbudget}, so the fourth cause fits its budget.

\subsection{Ordinary-line tags with a fixed moment order}
Use the line-cell allowance
\[
 e=\frac{1-\rho-\lambda}{4TC_{\rm line}}
 =\frac{9765625}{641885541199552},\qquad
 b_{\rm tag}=131458159,\qquad k_{\rm tag}=87,
\]
and fix \(m_{\rm tag}=b_{\rm tag}k_{\rm tag}+3=11436859836\).
Since \(b_{\rm tag}e/2\ge1\), Bernoulli's inequality implies
\[
 \left(\frac2{2+e}\right)^{m_{\rm tag}-3}\le2^{-87}.
\]
The principal term in \eqref{eq:tagprob} is bounded by
\begin{equation}\label{eq:fixedtagbudget}
 \frac{16C_{\rm line}}{e^2}\,2^{-87}
 =\frac{100590099609629337889425049}
 {219902325555200000000000000000000}
 <\frac1{2000000}.
\end{equation}
For this fixed order, \(m>2\) and \(1-(1-d)(m-1)<0\); hence
the remainder in \eqref{eq:tagremainder} tends to zero. Increase
\(n\) until \(\Theta_m(n)\le2\) and
\(C_{\rm line}R_m(n;e)\le1/2000000\). The tag probability is then
less than
\[
 2\left[\frac{16C_{\rm line}}{e^2}2^{-87}
          +\frac1{2000000}\right]<2\cdot10^{-6}.
\]
The credit identity and \eqref{eq:tagcapacity}, with this same
\(e\), ensure that a tag cannot hit its floor without a cell crossing.
This large moment order is one fixed integer; no expansion into
that many monomials is part of the numerical verification.

\subsection{The final rational comparison}
There are five causes, each with probability less than \(2\cdot10^{-6}\)
under the appropriate root point law. The maximum of their finitely
many root-uniform thresholds gives
\begin{equation}\label{eq:fixedfivebudget}
 \sup_p\Prob^{(p)}(E_p)<10^{-5}
\end{equation}
for every sufficiently large \(n\). The causes need not be independent
or disjoint. This is a bound for each root under its own point law,
not a simultaneous-success probability for all roots under one law.

\begin{corollary}[The numerical lower bound]\label{cor:target}
For every sufficiently large integer \(n\), there is a set
\(S\subseteq\{0,\ldots,n-1\}^2\) with no three distinct collinear
points and \(|S|\ge\lceil1501n/1000\rceil\).
\end{corollary}
\begin{proof}
The finite input satisfies Proposition~\ref{prop:contract}, with the
values above. The lattice comparison and positive margins supply
the hypotheses of the stopped process. Equation~\eqref{eq:massgrid}
gives \(A_n\ge A(1-512/n)^2\) for every integer \(n\ge512\).
By Proposition~\ref{prop:loss} and \eqref{eq:fixedfivebudget},
\(\E|S_n|/n\ge(1-10^{-5})A_n\). For \(n\ge2^{24}\), in
addition to the probability thresholds,
\begin{align}
 \frac{\E|S_n|}{n}-\frac{1501}{1000}
 &\ge(1-10^{-5})A(1-2^{-15})^2-\frac{1501}{1000}\notag\\
 &=\frac{2165079963272809251797}
 {274877906944000000000000000}>0.
 \label{eq:exactfinalcomparison}
\end{align}
Every output is legal by Proposition~\ref{prop:construction}.
Some output therefore has cardinality at least \(\lceil1501n/1000\rceil\).
All comparisons apply to every sufficiently large integer \(n\),
with no parity or divisibility restriction.
\end{proof}

A legal set has at most two points per row, so \(|S_n|\le2n\).
Equation~\eqref{eq:exactfinalcomparison} also gives
\[
 \Prob\{|S_n|\ge\lceil1501n/1000\rceil\}>\frac1{64000}
\]
past the same thresholds: on the complementary event \(|S_n|<1501n/1000\),
and the rational margin divided by \(2-1501/1000\) exceeds \(1/64000\).
The proof does not compute the final \(n_0\) or provide a practical
bound for small-grid simulations.

\section{Verification scope and provenance}\label{sec:status}

\subsection{Finite computation and mathematical proof}
The input consists of 1,024 nonnegative field integers and 262,144 ranks.
The latter assign all 242,176 positive type/cells to 481 phases.
Source, data and replay records are distributed under \path{verification/}.
The checker evaluates all four direction families directly, without
assuming reflection or type symmetry. The 481 freshly accepted records
use 385 tube bounds, 95 Farey bounds, and one global tube bound.

\noindent\textbf{Fresh execution for this edition.}
On 10 November 2026 the checker was rebuilt using GCC 13.3.0 on
Ubuntu under WSL, with four OpenMP threads and undefined-behavior
sanitization enabled. No cached phase result was reused.
All 481 phases and 242,176 positive type/cells passed; coverage checking
found no missing, duplicate or failing cell. The complete angular
replay took 1,355.4 seconds (about 22.6 minutes).
All 481 accepted CSV files are byte-identical to the incoming dossier's
records, and the maximum is exactly \eqref{eq:recordedrisk481}.
No sanitizer error was reported. Capacity and clock-constant checks
also passed. The separate arithmetic ledger passed 69 checks,
including actual field/rank reconstruction and agreement with this run.

The reproduction driver \path{verification/reproduce.py} requires a new
output directory and records the compiler, flags, input and source hashes,
executable hash, commands, exit codes and result hashes. Unsafe cached
resumption has been removed from the Python wrapper. The mathematical
C++ sources and fixed input are unchanged from the reviewed independent
implementation. Execution provenance and all per-phase results are
included in \path{verification/evidence/}; the machine-specific executable
is rebuilt from source rather than distributed.

The soundness arguments in Section~\ref{sec:certificates} connect finite
integer computations to geometric bounds. Hashes establish which files
were used; hashes alone do not prove those bounds. Likewise, exact
arithmetic certifies finite rational inequalities and field/rank
reconstruction, not the stochastic lemmas. Those lemmas are proved
in the preceding sections as natural-language mathematics.

The proof uses fixed orders and a fixed error budget. It needs neither
growing-order moment estimates nor an explicit value of the final \(n_0\).
The available partial Lean development concerns earlier parts of the
project and does not formally certify this manuscript or its 481-phase
application. No independent human review or complete formal verification
is claimed.

\subsection{Companion files}
The editable manuscript is \path{main.tex} and its section files.
\path{data/parameters.json} gives exact parameters and finite-sieve factors.
\path{scripts/check_arithmetic.py} reconstructs the input mass, phase
durations and supports from the field and rank files, then checks
capacity margins, discrete-risk slack, fixed-order witnesses and the
final rational budget. Its report is \path{evidence/arithmetic.json}.
Geometric replay, capacity checking and execution provenance have
separate reports in \path{verification/evidence/}; the arithmetic
script does not replace them. Source identifiers and hashes are
retained in the companion provenance records.

\subsection{AI use and licensing}
This work involved substantial use of OpenAI ChatGPT/Codex and AI reviews
supplied by the project maintainer. AI systems contributed mathematical
proposals and derivations, numerical programs, exact arithmetic checks,
Lean definitions and proofs, build orchestration, review, translation
and exposition. Separate AI contexts reviewed parts of the work. This
English edition was assembled and checked by AI assistance; separate
AI reviews are not independent human peer review.

The human maintainer chose the direction and requested experiments and
revisions as a hobby project. This disclosure does not attribute every
mathematical or programming step to human authorship. No named author,
affiliation or personal contact address is supplied. The repository's
hosting account is public. The project material and this manuscript
are distributed under the MIT license in \path{LICENSE}; retain its
notices when redistributing.

\begin{thebibliography}{9}
\bibitem{preparation}
Project preparation notes, \emph{Integrated proof draft preceding
Lean formalization}, 8--9 November 2026 (Japanese).
\href{https://github.com/agolkjaoi/no-three-in-line-lean/blob/282b636c160d049e3ef4d82f619fd44d98f8cf1d/docs/research-notes/preparation/PROOF_FOR_LEAN_JA.txt}{Fixed-version text}.
Historical source for the process and stopping arguments.
\bibitem{followup}
Project first followup, 9 November 2026 (Japanese).
\href{https://github.com/agolkjaoi/no-three-in-line-lean/blob/282b636c160d049e3ef4d82f619fd44d98f8cf1d/docs/research-notes/first-followup/FOLLOWUP_FOR_REVIEW_JA.txt}{Fixed-version text}.
Historical source for remaining exposure and the fixed loss budget.
\bibitem{second}
Project second followup, 9 November 2026 (Japanese).
\href{https://github.com/agolkjaoi/no-three-in-line-lean/blob/282b636c160d049e3ef4d82f619fd44d98f8cf1d/docs/research-notes/second-followup/SECOND_FOLLOWUP_FOR_REVIEW_JA.txt}{Fixed-version text}.
Historical source for the finite sieve and linewise cell count.
\bibitem{formal}
\emph{no-three-in-line-lean}, partial formalization snapshot, 10 November 2026.
\href{https://github.com/agolkjaoi/no-three-in-line-lean/tree/282b636c160d049e3ef4d82f619fd44d98f8cf1d}{Source snapshot}.
This is a partial formalization, not a formal certificate of the present
numerical theorem.
\bibitem{dossier}
Project 481-phase dossier and fixed-budget closure, 10 November 2026.
Finite input and executable verification sources are included in this
edition. The relevant arguments are developed in English here; the
historical Japanese supplements are not required for the mathematical
argument.
\end{thebibliography}
